\section{Results}
\label{sec:results}

\subsection{Per-drug standardization is what makes the multi-task objective learnable}

Table~\ref{tab:target} compares the three targets across representations and head counts, with all other
settings fixed. On the subset of \NLearn{} learnable drugs the three targets are statistically
indistinguishable, which shows that the curve fit itself buys no accuracy over the dose average. The
difference appears only at the full \NDrugs-head scale: both unstandardized targets collapse to
approximately zero, with scGPT turning negative, whereas the standardized target retains
$\rho=\rhoFullScgpt{}$ (scGPT). Per-drug standardization is therefore a prerequisite for a learnable
large multi-task objective rather than a change to the information content of the label.

\begin{table}[t]
\centering
\caption{Out-of-fold Spearman correlation, averaged over the \NLearn{} learnable drugs, seed 42. The three
targets tie at K\,$=$\,\NLearn{}, but only the standardized \code{auc\_z} survives at K\,$=$\,\NDrugs{}.
Values are means of the per-drug rows in \code{outputs/target/target\_comparison.csv}
(\code{notebooks/11}).}
\label{tab:target}
\begin{tabular}{llccc}
\toprule
Heads & Representation & \code{mean\_pv} & \code{auc} & \code{auc\_z} \\
\midrule
\NLearn{} & PCA   & 0.370 & 0.360 & \rhoTenPca{} \\
\NLearn{} & scGPT & 0.396 & 0.405 & \rhoTenScgpt{} \\
\NDrugs{} & PCA   & $+0.073$ & $-0.012$ & \rhoFullPca{} \\
\NDrugs{} & scGPT & $-0.078$ & $-0.069$ & \rhoFullScgpt{} \\
\bottomrule
\end{tabular}
\end{table}

A set of single-parameter interventions on the unstandardized \NDrugs-head configuration locates the cause
(Figure~\ref{fig:rescue}). Starting from the broken baseline ($\rho=-0.08$), heavier regularization, a
smaller model, a different batch size and sample reweighting all leave the result unchanged or make it
worse. Two changes help. Removing the regularization recovers most of the gap (to $\rho=+0.21$), which
shows that the dropout had been interacting with the mis-scaled loss by starving the learnable heads of
shared capacity. Standardizing the target recovers the signal fully (to $\rho=+0.33$) without weakening the
regularization, which is why, once the target is corrected, regularization is flat and no longer a useful
lever.

\begin{figure}[t]
\centering
\includegraphics[width=0.78\textwidth]{fig_rescue.png}
\caption{\textbf{On the broken \NDrugs-head configuration, only removing regularization or standardizing
the target recovers the signal.} Each bar applies a single change to the same unstandardized setting (raw
\code{auc}, scGPT). Heavier regularization, a smaller model, a different batch size and reweighting (grey)
do not help; removing regularization (orange) recovers most of the gap; per-drug standardization (blue)
recovers it fully and addresses the cause. Source:
\code{outputs/ablations/rescue\_k545.png} (\code{notebooks/10}).}
\label{fig:rescue}
\end{figure}

\subsection{The corrected model ranks unseen cell lines at full head count}

With the standardized target the model ranks unseen cell lines well above chance. On the learnable subset
it reaches $\rho=\rhoTenScgpt{}$ (scGPT) and $\rho=\rhoTenPca{}$ (PCA), and it holds at
$\rho=\rhoFullScgpt{}$ and $\rho=\rhoFullPca{}$ even when all \NDrugs{} heads are trained jointly, which
shows that the large multi-task objective is trainable once the target is standardized.

\subsection{Model capacity is exhausted; the ceiling is the number of independent cell lines}

On the corrected setting, model-side tuning yields nothing further: regularization, capacity (from
\ParamsHi{} down to \ParamsLo{} parameters) and batch size all leave the correlation flat. The decisive
control is a ridge regression fitted to the \NIndep{} cell-line mean embeddings, with no single cells and
no network, which reaches $\rho=\ridgePca{}$ (PCA) and essentially ties the PCA multi-layer perceptron
($\rho=\mlpPcaR{}$). The deep single-cell apparatus adds a measurable margin only through scGPT
($\rho=\mlpScgptR{}$), and even there the scGPT signal is nonlinear, since a linear head on scGPT drops to
$\rho=\linScgptR{}$. The scGPT representation does, however, generalize better than PCA in the way the
guiding hypothesis predicts: sharing an identical trunk, its training-to-validation gap on the corrected
target is \gapScgpt{}, against \gapPca{} for PCA, so the foundation-model embedding overfits far less. The
limiting factor is thus the small number of independent cell lines and the noise of the broadcast labels,
not the representation or the architecture.

\subsection{Against the field, the model clears the naive baseline but not the best-in-class model}

Under the DrEval leave-cell-line-out protocol and its normalized metric (Table~\ref{tab:dreval},
Figure~\ref{fig:dreval}), our model exceeds the NaiveMeanEffects baseline that half of the published field
fails to beat, reaching a normalized Spearman correlation of \drScgptSp{} for scGPT against 0.000 for the
baseline, and it slightly exceeds DrEval's own SingleDrug reference models on identical features
($\rho=\drRfSp{}$ for their random forest). It remains below the best model the DrEval study reports, whose
normalized coefficient of determination reaches about \drPaperRsq{} against our \drScgptRsq{}. The
scGPT-over-PCA advantage seen in the internal evaluation shrinks to near noise once the mean effects are
removed (normalized Spearman \drScgptSp{} versus \drPcaSp{}). Consistent with the ridge control above, this
reproduces the study's central finding that most of the apparent performance in this field is the
cell-line and drug mean effect rather than differential response.

\begin{table}[t]
\centering
\caption{DrEval leave-cell-line-out benchmark, normalized metrics, mean over five folds. The model clears
the NaiveMeanEffects baseline and edges DrEval's reference models on identical features. Data source:
\code{outputs/dreval/dreval\_lco\_results.csv} (\code{notebooks/12}).}
\label{tab:dreval}
\begin{tabular}{lcc}
\toprule
Model & norm.\ Spearman & norm.\ $R^2$ \\
\midrule
NaiveMeanEffects (baseline) & 0.000 & 0.000 \\
SingleDrugRandomForest (scGPT) & \drRfSp{} & 0.098 \\
OncoMLP (X\_pca) & \drPcaSp{} & \drPcaRsq{} \\
OncoMLP (X\_scGPT) & \drScgptSp{} & \drScgptRsq{} \\
\bottomrule
\end{tabular}
\end{table}

\begin{figure}[t]
\centering
\includegraphics[width=0.9\textwidth]{fig_dreval.png}
\caption{\textbf{The naive predictors' raw performance is bias that the normalization removes.} DrEval
leave-cell-line-out results, raw pooled correlation (left) and normalized (right). The naive predictors sit
above zero on the left and collapse to zero on the right, which shows that their raw performance was the
mean effect; the model remains above the baseline after normalization. Source:
\code{outputs/dreval/dreval\_lco.png} (\code{notebooks/12}).}
\label{fig:dreval}
\end{figure}
